\documentclass[12pt]{article}

\usepackage[utf8]{inputenc}
\usepackage[T1]{fontenc}
\usepackage[brazilian]{babel}
\usepackage{amsmath}
\usepackage{amssymb}
\usepackage{amsthm}
\usepackage{geometry}

\geometry{margin=2.5cm}

\newtheorem{proposicao}{Proposição}[subsection]
\newtheorem{afirmacao}[proposicao]{Afirmação}
\newtheorem{conjectura}[proposicao]{Conjectura}
\newtheorem*{definicao}{Definição}
\newtheorem*{afirmacaofalsa}{Afirmação Falsa}

\title{O que são provas matemáticas \\ \large (Lehman et al., 2018, pp. 3--8)}
\author{}
\date{}

\begin{document}

\maketitle

O objetivo principal desse texto é ensinar como utilizar do pensamento
matemático e metodológico para resolver e entender problemas computacionais.
Iniciaremos a nossa discussão entendo o que são as \emph{provas} nesse contexto
matemático, e também para o desenvolvimento desse pensamento analítico.

\begin{definicao}
	Uma \emph{prova matemática} de uma \emph{proposição} é uma série de
	\emph{deduções lógicas} que levam a uma proposição com base em uma corrente
	de \emph{axiomas}.
\end{definicao}

\section{O que é uma \emph{prova}?}

\subsection{Proposições}

\begin{definicao}
	Uma proposição é uma afirmação que pode ser tanto verdadeira como falsa.
\end{definicao}

Por exemplo, as duas afirmações abaixo são proposições. A primeira é
verdadeira, e a segunda é falsa.

\begin{proposicao}
	$1 + 2 = 3$
\end{proposicao}

\begin{proposicao}
	$1 + 1 = 3$
\end{proposicao}

Definir se uma proposição é falsa ou verdadeira nem sempre é simples. Na
maioria das vezes será necessário um estudo e uma análise por trás dela para
verificar se ela é verdadeira ou não, e é nesse momento que as \emph{provas}
são necessárias. Por exemplo:

\begin{afirmacao}
	Para todo inteiro $n$ não negativo, o valor de $n^{2} + n + 41$ é primo.
\end{afirmacao}

(Um número primo é um inteiro maior que 1 que não é divisível por nenhum outro
inteiro maior que 1. Por exemplo, 2, 3, 5, 7 e 11 são os cinco primeiros
números primos.)

Vamos tentar alguns experimentos numéricos para verificar essa proposição. Seja

\begin{equation}
	p(n) ::= n^{2} + n + 41
\end{equation}

Começamos com $p(0) = 41$, que é primo; então

$$
	p(1) = 43,\quad p(2) = 47,\quad p(3) = 53,\ \dots\ p(20) = 461
$$

são primos. Dessa maneira, essa afirmação começa a parecer plausível. Nós
conseguimos continuar testando até $n = 39$ e todos serão primos, mas se
fizermos $n = 40$ veremos que o resultado não é primo, pois

$$
	p(40) = 40^{2} + 40 + 41 = 41 \cdot 41
$$

que não é primo.

Dessa forma, a Afirmação 1.1.3 é falsa, já que $p(n)$ não é primo para todos os
$n$ inteiros não negativos.

Esse exemplo mostra algo pra nós. É impossível testar uma afirmação sobre
\emph{todos os números} existentes dentro de uma classe (que são infinitos)
testando apenas com um conjunto finito de números, não importa quão grande seja
esse conjunto.

A propósito, existe uma notação específica para proposições sobre \emph{todos}
os números de um determinado conjunto. Dessa maneira, a Afirmação 1.1.3 poderia
ser reescrita como

\begin{equation}
	\forall n \in \mathbb{N}.\ p(n)\ \text{é primo.}
\end{equation}

Aqui o símbolo $\forall$ é lido como ``para todos'', e o símbolo $\mathbb{N}$
diz respeito ao conjunto de todos os números inteiros não negativos $0, 1, 2,
	3, 4, \dots$ O símbolo $\in$ é lido como ``é membro de'' ou de maneira mais
curta, ``pertence a''. O ponto depois de $\mathbb{N}$ é apenas um separador de
frases.

Aqui está mais um exemplo:

\begin{conjectura}[Euler]
	A equação
	$$a^4 + b^4 + c^4 = d^4$$
	não tem solução quando $a, b, c, d$ são inteiros positivos.
\end{conjectura}

Euler (pronunciado ``oiler'') conjecturou isso em 1769. Mas a conjectura foi
provada falsa 218 anos depois por Noam Elkies em uma faculdade de artes
liberais na Mass Ave. A solução que ele encontrou foi $a = 95800,\ b = 217519,\
	c = 414560,\ d = 422481$.

Em notação lógica, a Conjectura de Euler poderia ser escrita,

$$
	\forall a \in \mathbb{Z}^+ \ \forall b \in \mathbb{Z}^+ \ \forall c \in
	\mathbb{Z}^+ \ \forall d \in \mathbb{Z}^+ .\ a^4 + b^4 + c^4 \neq d^4.
$$

Aqui, $\mathbb{Z}^+$ é um símbolo para os inteiros positivos. Sequências de
$\forall$'s como essa são geralmente abreviadas para facilitar a leitura:

$$
	\forall a, b, c, d \in \mathbb{Z}^+ .\ a^4 + b^4 + c^4 \neq d^4.
$$

Aqui está outra afirmação que seria difícil de refutar por amostragem: os
menores possíveis $x, y, z$ que satisfazem a igualdade têm cada um mais de 1000
dígitos.

\begin{afirmacaofalsa}
	$313(x^3 + y^3) = z^3$ não tem solução quando $x, y, z \in \mathbb{Z}^+$.
\end{afirmacaofalsa}

Vale a pena mencionar mais um par de proposições famosas cujas provas foram
buscadas por séculos antes de finalmente serem descobertas:

\begin{proposicao}[Teorema das Quatro Cores]
	Todo mapa pode ser colorido com 4 cores de modo que regiões adjacentes tenham
	cores diferentes.
\end{proposicao}

Várias provas incorretas desse teorema foram publicadas, incluindo uma que
resistiu por 10 anos no final do século XIX antes que seu erro fosse
encontrado. Uma prova trabalhosa foi finalmente encontrada em 1976 pelos
matemáticos Appel e Haken, que usaram um programa de computador complexo para
categorizar os mapas coloríveis com quatro cores. O programa deixou alguns
milhares de mapas sem categorizar, os quais foram verificados manualmente por
Haken e seus assistentes --- entre eles sua filha de 15 anos.

Havia motivo para duvidar se essa era uma prova legítima --- a prova era grande
demais para ser verificada sem um computador. Ninguém podia garantir que o
computador havia calculado corretamente, nem havia entusiasmo em despender o
esforço de reverificar manualmente as colorações de quatro cores de milhares de
mapas. Duas décadas depois, foi encontrada uma prova majoritariamente
inteligível do Teorema das Quatro Cores, embora um computador ainda seja
necessário para verificar a colorabilidade com quatro cores de várias centenas
de mapas especiais.

\begin{proposicao}[Último Teorema de Fermat]
	Não existem inteiros positivos $x$, $y$ e $z$ tais que
	$$x^n + y^n = z^n$$
	para algum inteiro $n > 2$.
\end{proposicao}

Em um livro que estava lendo por volta de 1630, Fermat afirmou ter uma prova
para essa proposição, mas não havia espaço suficiente na margem para
escrevê-la. Ao longo dos anos, o Teorema foi provado válido para todo $n$ até
4.000.000, mas já vimos que isso não deveria necessariamente inspirar confiança
de que ele vale para todo $n$. Há, afinal, uma clara semelhança entre o Último
Teorema de Fermat e a falsa Conjectura de Euler. Finalmente, em 1994, o
matemático britânico Andrew Wiles apresentou uma prova, após sete anos
trabalhando em segredo e isolamento em seu sótão. Sua prova não cabia em
nenhuma margem.

Por fim, vamos mencionar mais uma proposição simplesmente enunciada cuja
veracidade permanece desconhecida.

\begin{conjectura}[Goldbach]
	Todo inteiro par maior que 2 é a soma de dois números primos.
\end{conjectura}

A Conjectura de Goldbach remonta a 1742. Sabe-se que ela vale para todos os
números até $10^{18}$, mas até hoje ninguém sabe se ela é verdadeira ou falsa.

Para um cientista da computação, algumas das coisas mais importantes a se
provar são a corretude de programas e sistemas, se um programa ou sistema faz o
que deveria fazer. Programas são notoriamente cheios de erros, e há uma
comunidade crescente de pesquisadores e profissionais tentando encontrar formas
de provar a corretude de programas. Esses esforços têm sido bem-sucedidos o
suficiente no caso de chips de CPU a ponto de serem rotineiramente usados pelos
principais fabricantes de chips para provar a corretude dos chips e evitar
alguns erros notórios do passado.

Desenvolver métodos matemáticos para verificar programas e sistemas continua
sendo uma área de pesquisa ativa. Ilustraremos alguns desses métodos no
Capítulo 5.

\end{document}
